\subsection{Ordered fields}

DEFINITION 2. --- A commutative field, equipped with a total ordering, is called an ordered field if its ordering and its ring structure are compatible.

We restrict ourselves to total order relations on fields because the others are very « pathological » (VI, p.~38, Ex. 6).

Examples. --- 1) The field \(\mathbf{Q}\) of rational numbers is an ordered field.

\begin{enumerate}
\def\labelenumi{\arabic{enumi})}
\setcounter{enumi}{1}
\item
  A subfield of an ordered field, with the induced ordering, is an ordered field.
\item
  * The field of real numbers is an ordered field. *
\end{enumerate}

Let \(\mathbf{K}\) be an ordered field. For all \(x\in \mathbf{K}\) we put

\[
\begin{array}{l l} \operatorname{sgn} (x) = 1 & \text {if} \quad x > 0, \\ \operatorname{sgn} (x) = - 1 & \text {if} \quad x <   0, \\ \operatorname{sgn} (x) = 0 & \text {if} \quad x = 0. \end{array}
\]

Then we have \(\operatorname{sgn}(xy) = \operatorname{sgn}(x)\operatorname{sgn}(y)\) ; we call \(\operatorname{sgn}(x)\) the sign of \(x\) . The map \(x \mapsto \operatorname{sgn}(x)\) from \(\mathbf{K}^*\) to the multiplicative group \(\{-1, +1\}\) is a surjective homomorphism whose kernel, the set of strictly positive elements of \(\mathbf{K}\) , is a subgroup of \(\mathbf{K}^*\) of index 2.

Conversely, if K is a commutative field and \(s: K^{*} \to \{-1, +1\}\) is a surjective homomorphism whose kernel is closed under addition, then s is the sign map for a unique ordered field structure, where the set of strictly positive elements is the kernel of s.

For all \(x\) and \(y\) in \(\mathbf{K}\) we have \(x = \operatorname{sgn}(x)|x|\) and \(|xy| = |x||y|\) .

On the other hand every ordered field is of characteristic zero (Prop. 1).

PROPOSITION 2. --- Let A be a totally ordered integral domain, and let K be its field of fractions. Then there exists one and only one ordering on K which restricts to the given ordering on A and makes K an ordered field.

Every \(x \in K\) can be expressed in the form \(x = ab^{-1}\) , with a and b in A and \(b \neq 0\) . If x is positive, then a and b have the same sign, and conversely. Thus we see that, if there exists an ordering on K satisfying the prescribed conditions, then it is unique, and the set P of positive elements is identical to the set of \(ab^{-1}\) , where a and b are elements of A of the same sign, and \(b \neq 0\) . It remains to show that P satisfies conditions \((\mathrm{AP}_{\mathrm{I}}), (\mathrm{AP}_{\mathrm{II}}), (\mathrm{AP}_{\mathrm{III}})\) and \((\mathrm{AP}_{\mathrm{IV}})\) . This is obvious for \((\mathrm{AP}_{\mathrm{II}})\) and \((\mathrm{AP}_{\mathrm{IV}})\) . For \((\mathrm{AP}_{\mathrm{I}})\) , consider \(ab^{-1} + cd^{-1}\) , where we may assume that \(a, b, c\) and \(d\) are positive; this sum is \((ad + bc)(bd)^{-1}\) , and \(ad + bc\) and \(bd\) are positive.

To show \((\mathrm{AP}_{\mathrm{III}})\) , consider an identity of the form \(ab^{-1} = -cd^{-1}\) , so that \(ad + bc = 0\) . If we assume that a and b have the same sign and that c and d have the same sign, then the sign rules show that ad and bc have the same sign; whence ad = bc = 0, so a = c = 0; hence P does indeed satisfy \((\mathrm{AP}_{\mathrm{III}})\) .

Example. --- Since Z admits only one totally ordered ring structure (VI, p.~19, example), the field Q admits only one ordering which makes it an ordered field: this is the usual ordering.

